% Fragmento público generado desde propietarios íntegros.
% Residencia: 52.1.2.
% Fuente: ../incoming/radion_monografia_20260827/manuscrito/sections/02_genealogia.tex:181-272

\subsubsection{Genealogía de las constantes y construcción del par angular}

\paragraph{Sección coinductiva y unicidad del valor generado de \(\pi\)}

Sea \(x_\infty\) la sección global del sistema inverso. Los lectores
coinductivos producen
\[
\operatorname{ev}_\pi(x_\infty)=\pih,\qquad
\operatorname{ev}_e(x_\infty)=e_{\HMT},\qquad
\operatorname{ev}_\varphi(x_\infty)=\varphi_{\HMT}.
\]
El subíndice HMT no designa otro número. Conserva la genealogía de la misma
constante. En particular,
\[
\pih=3.14159265358979323846264338327950288419716939937510\ldots.
\]
La vuelta positiva es
\begin{equation}
T_+(x_\infty)=2\pih.
\label{eq:vuelta-plus}
\end{equation}
Su publicación angular de una unidad es
\begin{equation}
R_{\HMT}^{\circ}=\frac{360^\circ}{T_+}=\frac{180^\circ}{\pih}.
\label{eq:radian-publicado}
\end{equation}
La ecuación no define todavía el radión enriquecido; define la carta visible
del radián una vez generada la vuelta.

\paragraph{Bidireccionalidad de rutas y prolongaciones conjugadas}

La bidireccionalidad no se enuncia como eslogan. La involución concreta
\begin{equation}
J_\alpha(T)=\frac{\alphah^2}{T},\qquad
T_-=J_\alpha(T_+),\qquad T_+T_-=\alphah^2
\label{eq:bidireccional-turn}
\end{equation}
conserva un producto y revierte la hoja. Expansión y contracción son dos
prolongaciones conjugadas del mismo estado. La raíz finita de un jet
cuadrático puede aproximar \(T_+\), pero no sustituye a la sección
coinductiva: la diferencia certificada es
\[
T_{P_9}-T_+=-5.1110566290335\ldots\times10^{-25}.
\]
El control finito es un testigo; el fundamento es el límite.

\paragraph{Dependencia causal entre \(\alpha\), acción y publicación angular}

El cierre dodecafásico genera \(\alphah\) desde el estado firmado y sus
cartas reversibles. Después, el lector determinantal local fija la década de
acción mediante
\[
\Sigma_H=\varphi_{\HMT}\alphah^{16},
\qquad \operatorname{ord}_{10}\Sigma_H=-34.
\]
La sección de retorno produce \(\hret\). Ni \(\alpha\) ni \(\hbar\) se
\enquote{miden en grados}. Se publican posteriormente mediante el par angular
\begin{equation}
P(\alpha,\eta_{\mathrm{ret}})=
\left([10^3\alpha]_{360},[2(\eta_{\mathrm{ret}}+\alpha)]_{360}\right)
=(A,C^*).
\label{eq:parangular}
\end{equation}
En la rama canónica,
\[
A=7.297352569283800997285105472380663\ldots{}^\circ,
\]
\[
C^*=2.1237383389686457242619513803933607937\ldots{}^\circ.
\]
Sólo entonces la carta analítica
\[
x=\frac{\pi A}{180},\qquad y=\frac{\pi C^*}{180}
\]
abre los canales \(q_\pm=e^{-x\mp y}\).

El orden causal completo es
\[
\alpha\to(\eta\parallel-34)\to\hbar
\to P\to\kappa_{360}\to q_\pm
\to\text{realizaciones MD}.
\]
La relación \(C^*/2-\alpha\) puede servir de verificación inversa. No se usa
para generar \(\hret\) retroactivamente.

\begin{narrativebox}
La carta angular no convierte una constante adimensional o una acción en un
ángulo por decreto. Publica dos coordenadas del estado en una rueda finita.
El grado es el idioma de esa publicación; no es la sustancia ontológica de
\(\alpha\) ni de \(\hbar\).
\end{narrativebox}
% Fuente: ../incoming/radion_monografia_20260827/manuscrito/sections/03_cierre_exacto.tex:254-285
\paragraph{Independencia prospectiva respecto del valor final}

El verificador del operador de profundidad prohíbe expresamente tres entradas:
\[
\frac{180}{\pi},\qquad \epsR,\qquad
\text{el decimal objetivo}.
\]
Primero construye \(S_E\) y \(S_T\), después ejecuta la resta APP, y sólo al
final compara el cierre. Los controles dan
\[
\texttt{target\_epsilon\_used=False},\qquad
\texttt{target\_180\_over\_pi\_used=False}.
\]
Asimismo:
\begin{itemize}
\item una sola hoja da un defecto de orden \(10^{-5}\);
\item sustituir \(\theta_2\) por \(\theta_1\) o \(\theta_3\) desplaza el
resultado en aproximadamente \(\pm\pi/6\);
\item eliminar el canal \(e\) deja un defecto de orden \(10^{-3}\);
\item la variante histórica \(\Delta_4(A,C^*)\) produce otro número;
\item la cola espectral posterior a \(\Rstate\) y el operador armónico de
cuatro términos tampoco coinciden con \(\epsone\).
\end{itemize}

\begin{statusbox}
\textbf{Estatuto.} El cierre es una \status{formalización nueva} demostrada
con estructura de partida explícita y certificado computacional independiente del valor final.
La igualdad es exacta en el límite coinductivo. El perfil corto de 36 cifras
de \(\alpha\) genera una variante que diverge sólo después de muchas cifras;
se conserva como testigo histórico, no como valor canónico.
\end{statusbox}

